\begin{resumo}[Abstract]
\begin{otherlanguage*}{english}
This technical report documents the development of the \textbf{Karaokê Apache} system, a university extension project carried out for the \textbf{Apache Gru Rock Bar}, in Guarulhos, São Paulo, Brazil. During the venue's Karaoke Night, song requests were made verbally or on paper, and the host searched for each video manually on YouTube through the projector, which led to lost requests, repeated songs and a lack of visibility over the singing order. The goal of the project is to develop the Minimum Viable Product (MVP) of a low-cost system to receive requests, organize the singer queue and play the videos on the main screen. Based on requirements gathered with the partner, five business rules were defined, including first-come, first-served ordering (FIFO), blocking of repeated songs and a limit of requests per customer, and the system was modeled with UML. The architecture comprises React web interfaces for customers and the host, accessed through a QR Code shown on the main screen; a Java 21 API built with Spring Boot and a MySQL database, both hosted in the cloud, which searches videos through the \textit{YouTube Data API v3}; and an Electron player, running on the bar's computer, which shows the current song's video next to the system's own interface on the main screen. Queue changes are pushed to all screens over WebSocket using the STOMP protocol. The Proof of Concept validated real-time synchronization between the web interface and the player, and the MVP is under implementation, with tests planned for the business rules, access control and update time. The only expected cost for the venue is cloud hosting.

\vspace{\onelineskip}
\noindent
\textbf{Keywords}: Karaoke; Queue management; Web application; Electron; University extension.
\end{otherlanguage*}
\end{resumo}
